\chapter{Performance}\label{ch:performance}

The measured resource model of one G17 core has six resources, and a kernel's time is set by whichever it
exhausts first:

\begin{enumerate}
\def\labelenumi{\arabic{enumi}.}
\item Instruction issue: one simdgroup issues a simple instruction about every 3.1~ns; there is no fp32/int32 dual
  issue (\cref{sec:scalar-issue}).
\item The tensor unit: about 4.96~ns per 16 × 16 × 16 MMA per core at saturation, 1,024 flop per cycle per core
  (\cref{sec:tensor-issue-rate}).
\item Dependent latency: a dependent MMA or a load-to-use chain in one simdgroup costs far more than the issue
  interval (74~ns per MMA in one chain, 30.5~ns per L1 load, 405–415~ns per DRAM load); only independent work, from
  the same simdgroup or from other resident simdgroups, hides it (\cref{sec:tensor-issue-rate,sec:memory-bandwidth-footprint}).
\item Residency: how many simdgroups a core holds, which is what supplies that independent work (\cref{sec:residency-occupancy-measurements}).
\item Memory bandwidth by footprint: on-core, SLC (520–528~GB/s) and DRAM (269.6–298.5~GB/s depending on the
  instrument) (\cref{sec:memory-bandwidth-footprint}).
\item Instruction fetch and dispatch: a hot loop body over 16 KiB loses the instruction cache at low occupancy; each
  dependent dispatch costs microseconds (\cref{sec:memory-bandwidth-footprint,sec:dispatch-encoder-command-buffer}).
\end{enumerate}

In the kernels this project studied, the tensor unit's peak was rarely the limit. The fused attention and FFN kernels execute about 46 instructions per MMA, so the tensor pipe sits
idle most of the time and they are issue- and ALU-bound at 35–57 percent of peak (\cref{sec:measured-bottlenecks-studied}). The other
studied workloads are bound elsewhere: the one-token decode projections are DRAM-bound at q8 and
issue-bound by counter at q4 (yet the q4 MLP did not speed up when a quarter of its instructions were removed);
batched decode is bound first by occupancy and, after split-K, by the fp32 pipe; prefill attention by latency on a
scratch round trip; and prefill GEMMs reach 80–89 percent of peak once measured in the graph (\cref{sec:bounds-real-workloads}).

Rate and timing terms:

\begin{itemize}
\item \textbf{Per core} / \textbf{per simdgroup}: a rate quoted per core aggregates every simdgroup on it; a per-simdgroup rate is
  one instruction stream. The device has 20 cores and four SIMD slots per core (\cref{ch:device-execution-model}).
\item \textbf{Occupancy} or \textbf{simdgroups per core}: simdgroups launched divided by 20 (the launched figure) or resident at once
  (the residency figure); the text says which.
\item \textbf{Saturated}: enough simdgroups per core that adding more does not raise throughput.
\item \textbf{Warm clock}: timed after a saturating dispatch long enough to raise the GPU to its top performance state
  (\cref{sec:clocks-power-gpu}). Every absolute time in this chapter is at a warm clock unless stated.
\item \textbf{Footprint}: unique bytes a kernel touches.
\item MMA opcode names are those of \cref{ch:tensor-unit}; the tensor MMA used for rate measurements is \op{5106}.
\end{itemize}

Beyond the 0.4–1.3 percent in-process and 9–13 percent cross-process spread of \cref{sec:clocks-performance-states},
one row differed by 18 percent between processes, and the rebuilt harness moved 25–44
percent between its own processes; ratios timed together in one process reproduce to about 1.5 percent (measured,
MM~\mm{16}, MM~\mm{25.124.2}), which is why the comparisons here are ratios. Under load the second run of a fixed A-then-B pair ran
about 0.1~ms per token faster, so every A/B here is order-balanced (measured, MM~\mm{17} rule 43).

\section{Tensor issue rate and latency}\label{sec:tensor-issue-rate}

At saturation (measured, MM~\mm{8}), 16,384 threadgroups x 4 simdgroups x 1,400 issues of the 16-bit MMA, \op{5106}, is 91.75 M issues in
22.44~ms: 4.09 × $10^{9}$ issues per second, 33.5 TFLOP/s dense fp16/bf16 into fp32 at 8,192 flop per issue, 1.67
TFLOP/s per core, 128 MAC per cycle per SIMD (four times a 32-lane FMA) at an implied 1.63~GHz. int8 runs at exactly
twice the rate: 13~ns per issue in a chain, 11.50~ms saturated, 8.1 × $10^{9}$ issues per second, 67 TOPS, 256 MAC per cycle
per SIMD, using \op{10384}.

An independent later harness gives 4.96–4.97~ns per MMA per core at 16 simdgroups per core, 8,192 flop in 8.0
cycles at 1.62~GHz, 1,024 flop per cycle per core; int8 2.47~ns. About 5.5 simdgroups per core saturate the unit
(measured, MM~\mm{8}, MM~\mm{0.1}). Percent-of-peak figures in the sources use 33.1 TFLOP/s, carried from that harness and not
re-derived; 33.5 is the derived figure. They differ by 1.2 percent, inside harness spread (MM~\mm{8}). This chapter quotes
percentages against 33.1, as the sources do.

Each core has four tensor slots (measured, MM~\mm{8}): one, two or four simdgroups each running the chain take the same
time (32.8~µs), 8 take 57.5 and 32 take 224.3, flat to four and linear beyond (the full row is in \cref{sec:core-logical-resource}). Each
SIMD has one slot and holds it for the whole interval.

The cost of one MMA depends on how much independent work surrounds it (\cref{tab:mma-latency-occupancy}).

\begin{xltabular}{\linewidth}{@{}>{\hsize=1.12\hsize}X>{\hsize=1.25\hsize}X>{\hsize=0.62\hsize}X@{}}
\caption{Cost per MMA by occupancy and by dependence between MMAs.}\label{tab:mma-latency-occupancy}\\
\toprule
Configuration & Cost per MMA & Source \\
\midrule
\endfirsthead
\tablecontinued{3}
\toprule
Configuration & Cost per MMA & Source \\
\midrule
\endhead
dependent chain at saturation & 22.2~ns; eight independent accumulators over the same 1,400 issues take
the same 31.96~µs, so result latency is no longer than the issue
interval & measured, MM~\mm{8} \\*
one simdgroup, 8 independent accumulators & 20.5~ns & measured, MM~\mm{8} \\*
one simdgroup, one dependent chain & 74~ns (3.6 times) & measured, MM~\mm{8} \\*
one simdgroup per core: 8 dependent MMAs vs four independent chains of 2 & both 218.0~ns per iteration, so chaining adds nothing and the cost is that
of a single chain & measured, MM~\mm{8} \\*
fit over 321 kernel timings, bf16 class & 62.6~ns inside a dependent chain at 1 simdgroup per core; 5.8~ns at 64 & measured fit, MM~\mm{16} \\*
same fit, fp8 and int8 & 68.2~ns at 1; 7.0~ns at 64 & measured fit, MM~\mm{16} \\*
\op{5107} back-to-back
into the same destination & 34~ns (a write-after-write stall), gone across 8 destinations or at
saturation & measured, MM~\mm{8} \\*
\bottomrule
\end{xltabular}

The same fit charges 1.6~ns per other instruction at one simdgroup and 0.21–0.34~ns at 64. Its six worst outliers are
all independent-MMA ablations, so the large per-MMA figure measures dependence rather than the MMA itself (MM~\mm{16}). The
fp32-operand forms (\cref{sec:instruction-family}) and a transposed A cost nothing extra (measured, MM~\mm{8}).

At low occupancy a compiler should interleave independent accumulators (one simdgroup goes from 74 to 20.5~ns
per MMA); at 12 or more simdgroups per core the forms converge and the interleave buys nothing (MM~\mm{8}, MM~\mm{11}). A
255-trip one-slice K loop in one simdgroup averages 1.45~µs per trip (loads then four MMAs, nothing overlapping), the
low-occupancy extreme (measured, MM~\mm{25.101}).

The static model's tensor constants are fitted, not measured (\cref{sec:cost-models-their}; MM~\mm{25.144.12}).

\textbf{Evidence.} MM~\mm{0.1}, MM~\mm{8}, MM~\mm{11}, MM~\mm{16}, MM~\mm{25.101}, MM~\mm{25.144.12}.

\section{Bottlenecks in the studied kernels}\label{sec:measured-bottlenecks-studied}

Fused attention and FFN are issue- and ALU-bound (measured, MM~\mm{8}, MM~\mm{0.12}, MM~\mm{16}). Fused attention reaches 35–42
percent of 33.1 TFLOP/s (11.5–13.8 TFLOP/s) and the FFN 47–57 percent. A real attention chunk loop executes about
46 instructions per MMA: 1,480 instructions for 32 MMAs is 356~ns of the 433~ns a key block costs per core, while
the tensor pipe is busy for 158~ns. Attention throughput stays at 11.5–13.8 TFLOP/s from an 8 MiB working set to 128
MiB, so it is issue-bound, not bound by the KV stream. One simdgroup's dependent chain costs
5.35~µs per 32-key block; about 247 simdgroups saturate the GPU.

In the attention kernels (measured, MM~\mm{16}, recon kernels), simdgroups per threadgroup from 1 to 16
are within noise (keep at least 20 threadgroups so every core has one); transposed-B flags cost 1.3 percent; unrolling
K by 2 or 4 gains nothing and adds 98–110 spill stores; two token tiles per task doubles the time (354 spill stores).
The retained rebuild costs 1.25–1.29 times the recon per key block (6.9 against 5.35~µs for one simdgroup; 558 against
433~ns per block and core) despite 1,033 instructions and no spills, against the recon's 1,925 with 42 (measured, MM
25.126).

Loop form and fusion were measured in MM~\mm{16}. The runtime K loop is 1.3–1.8 times faster than unrolling at 8
blocks and 5.9–7.2 times at 32 (unrolled kernels spill 24–923 stores). A spill-free fused chain is 5–10 percent faster than any cut and 8–20
percent faster than three kernels; at 15 or more spill stores the cuts win. Attention stays fused in all 27 measured
shapes even while spilling 42 stores (cutting P costs 1.03–1.30 times, cutting both 1.14–2.03 times). For a gated FFN
at d = 2048, Hh = 8192, one wave is 20 threadgroups at about 1.8~ms in fp8, and the fused form needs \code{nsg*mt*640}
bytes of threadgroup memory (fp8, int8) or \code{nsg*mt*1,088} (bf16, fp16) against 32,768: \code{nsg*mt <= 51} and \code{<= 30}.

\Cref{tab:gemm-vs-matmul2d} compares this compiler's GEMM with Apple's \code{matmul2d} (measured, interleaved medians, M × 128 × 256 half, MM~\mm{25.107}).

\begin{xltabular}{\linewidth}{@{}>{\hsize=1.15\hsize}Xl>{\hsize=0.85\hsize}Xl@{}}
\caption{GEMM time for M × 128 × 256 half: Apple's \code{matmul2d} against this compiler's K-loop and unrolled forms.}\label{tab:gemm-vs-matmul2d}\\
\toprule
Shape and grid & Apple & this compiler, K loop & unrolled \\
\midrule
\endfirsthead
\tablecontinued{4}
\toprule
Shape and grid & Apple & this compiler, K loop & unrolled \\
\midrule
\endhead
M512, 8 threadgroups & 106.6~µs & 113.7~µs (1.07× behind) & 262.7~µs \\*
M1024, 8 threadgroups & 183.6~µs & 244.7~µs (1.33×) & 376.0~µs \\*
M4096, 256 threadgroups x 1 simdgroup (equal geometry) & 58.7~µs & 62.0~µs (1.06×) & 78.6~µs \\*
\bottomrule
\end{xltabular}

Apple's own 64 × 4 geometry takes 69.4~µs for M4096; at 12.8 simdgroups per core the K loop reaches 4,330 GFLOP/s
against Apple's 4,570. \emph{Correction:} MM~\mm{25.98}'s earlier ratios are superseded by MM~\mm{25.107} (\cref{app:a}).

\begin{itemize}
\item Occupancy mattered more than instruction scheduling. An occupancy sweep (one 16-row tile row per simdgroup, a 12.5
  KB body) goes from 132 GFLOP/s at 8 simdgroups (0.4 per core) to 3,329 at 256 (12.8 per core) and 4,497 at 512 (25.6
  per core, 124 registers); work grew 32 times and time 1.27 times (measured, MM~\mm{25.98.1}). Software pipelining
  (rotating MMAs ahead of loads) gave 1.08 times at 8 simdgroups (113.8 to 105.6~µs, where at least 1.2 had been
  preregistered) and 0.99 times at 12.8 per core (measured, MM~\mm{25.97}, MM~\mm{25.107}).
\item In the K-loop body, one slice per trip waiting on its own loads paid a full round trip per step (about 417~ns against
  Apple's about 245 at 8 threadgroups). Issuing two slices' loads before the first MMA (\code{kloop\_unroll=2}) gains 35–38
  percent and closes to 3–4 percent of Apple at 4 and 8 threadgroups (about 337 to 210~ns per step against Apple's
  202); an interleaved control gains 1–4 percent (measured, MM~\mm{25.124.6}).
\item Unrolled bodies are fetch-bound, at about 518~ns per MMA (512 MMAs in 265~µs) against a 20–74~ns issue interval;
  fully unrolled K=256 bodies are 61~KB and 116~KB against Apple's 5,142-byte looping kernel (measured, MM~\mm{25.97},
  MM~\mm{25.98}).
\item Grid-split threadgroups run concurrently: with 512 MMAs per threadgroup, time is flat from G = 1 to 8
  (\code{t(8)/t(2) = 0.995}); the null control (2 MMAs) is 11.5 against 265.3~µs (measured, MM~\mm{25.89}).
\end{itemize}

For the bf16 layer kernels (measured, MM~\mm{25.43}), 16 simdgroups per threadgroup is fastest for the two- and
three-kernel forms (2.856 and 3.103~ms); 32 is 1.46–1.86 times slower than 16 in every form, against a run spread of
0.007–0.069~ms.

A weight-stationary FFN is 3.7–4.4 times faster than a token-parallel one at 16 and 32 tokens (measured, MM~\mm{16},
MM~\mm{25.126}); on the retained harness fp8 weight-stationary loses from 48 tokens (1.13 times the
token-parallel time) and bf16 from 64 (1.15). \emph{Correction:} the recon model's "about 64 tokens" crossover is not the
measured one for these kernels.

A GELU register epilogue (x * sigmoid(1.702×), six instructions per slot, 4 registers)
takes 15.6~µs at M16 against 52.2~µs for a memory stage (MM~\mm{25.109}). A full-tile one-dispatch attention takes 30.5~µs
at 2 rows to 104.3~µs at 16 rows, because its softmax is straight-line (MM~\mm{25.110}).

\textbf{Evidence.} MM~\mm{0.12}, MM~\mm{8}, MM~\mm{16}, MM~\mm{25.43}, MM~\mm{25.89}, MM~\mm{25.97}, MM~\mm{25.98}, MM~\mm{25.98.1}, MM~\mm{25.107}, MM~\mm{25.109},
MM~\mm{25.110}, MM~\mm{25.124.6}, MM~\mm{25.126}.

\section{Scalar issue}\label{sec:scalar-issue}

\Cref{tab:scalar-issue-pairs} gives the issue model (measured, MM~\mm{25.142.6}, MM~\mm{0.15}; a counted 128-trip loop over eight
independent chains, op counts checked from the decoded bytes, each dispatch in its own process).

\widetable
\begin{xltabular}{\linewidth}{@{}>{\hsize=1.53\hsize}Xlll>{\hsize=0.74\hsize}X>{\hsize=0.74\hsize}X@{}}
\caption{Scalar issue cost per op in one simdgroup, each class alone and mixed, against the predictions of one shared port and of dual issue.}\label{tab:scalar-issue-pairs}\\
\toprule
One simdgroup, ns per op & alone A & alone B & mixed & one shared port predicts & dual issue predicts \\
\midrule
\endfirsthead
\tablecontinued{6}
\toprule
One simdgroup, ns per op & alone A & alone B & mixed & one shared port predicts & dual issue predicts \\
\midrule
\endhead
ffma + iadd & 4.18 & 3.09 & \textbf{3.61} & 3.64 & 2.09 \\*
fadd + iadd & 3.11 & 3.09 & 3.11 & 3.10 & 1.55 \\*
fmul + shl\_imm & 3.16 & 6.10 & 3.17 & 4.63 & 3.05 \\*
fadd + fmul (control, one class) & 3.11 & 3.16 & 3.17 & 3.14 & 1.58 \\*
\bottomrule
\end{xltabular}

\begin{itemize}
\item A simple instruction costs about 3.1~ns per simdgroup whatever its class; this is the per-simdgroup issue floor.
\item There is no fp32/int32 dual issue. The ffma + iadd pair separates the two models, and it lands on the shared-port value. The full grid
  agrees: fadd 1,050, iadd 1,238, mixed 1,271~µs at B = 128.
\item fp32 ffma at 4.18 shows its latency is not fully hidden by eight chains.
\item \code{shl\_imm} costs 6.1 alone and 3.17 mixed with fmul: a stall that other instructions fill (the narrow register file's
  constraint, \cref{ch:register-file}), not an issue slot.
\item fp32 add throughput is about 4.1 × $10^{12}$ per second machine-wide, about 2.1 × $10^{11}$ per core, consistent with 128
  fp32 lanes per core at about 1.6~GHz (the clock was not measured by this probe).
\end{itemize}

For a compiler, interleaving float and integer work gains nothing; independent chains do, and eight reach
the floor (MM~\mm{17} rule 38). Apple's wait after every special-register read costs nothing measurable (\cref{sec:read-latency-whether},
MM~\mm{25.121}). The static model's fitted issue constants (3.105~ns interval, 12.5 issue slots per interval per core, 3.43
ns per fragment load, about 150~GB/s of fragments per core) are in \cref{sec:cost-models-their} (MM~\mm{25.144.12}).

ALU numerics edges that affect cost are in \cref{ch:scalar-arithmetic}. The one relevant to scheduling is that \code{fdiv} costs 9 more instructions
than a single-instruction operation such as \op{1272} (measured, MM~\mm{6}).

\textbf{Evidence.} MM~\mm{0.15}, MM~\mm{6}, MM~\mm{17}, MM~\mm{25.121}, MM~\mm{25.142.6}, MM~\mm{25.144.12}.

\section{Cost of registers and spills}\label{sec:registers-spills-occupancy}

The register-file and capacity rules (the 126-register flat pool, the trade lines, the accumulator limits, the two
spill opcodes, spill cliffs as scheduling outcomes) are in \cref{sec:capacity-budgets-spills}.

A spill costs latency rather than saturated throughput (measured, MM~\mm{11}). Against a 14-accumulator kernel, 15
accumulators (32 spill stores) cost 1.31 times, 16 cost 1.47, 18 cost 2.17 and 24 cost 2.30 at 1–4 simdgroups per
core, but only 1.00–1.10 times at 64. Removing spills did not always help; with few stores the streaming order ran
at 0.47–0.79 of the flat order's speed in several shapes.
Spill counts here include both spill forms (\cref{sec:capacity-budgets-spills}; MM~\mm{11}, MM~\mm{13}, MM~\mm{19}).

For this compiler's one-tile GEMM body (measured, MM~\mm{25.111}), a declared count
of 44 against 124 (one metadata byte) is equal in time within 1.2 percent at 25.6 simdgroups per core (41.6 against
42.1~µs) and 0.7 percent at 51.2 per core (69.9 against 69.4~µs). The result is specific to this body, and equal time
past 32 simdgroups per core does not by itself show residency past 32.
Register peaks of this compiler: 16×16×64 R42, 32×32×64 R82, 64×64×64 R122; M512 with GELU reaches R125; declared 44,
60, 92 and 124 at N16, 32, 64 and 128 (this compiler, MM~\mm{25.98.1}, MM~\mm{25.109.1}).

Registers did cost time in the decode kernels (measured, MM~\mm{25.144.4}). In the decode qmv kernel at 4 rows per simdgroup, the same
loop at 81 registers instead of 114 is 4.1~µs faster (43.2 against 47.3~µs). The B~8 batched qkv kernel uses 109 (q4)
and 108 (q8) registers and holds occupancy at 25–27 percent where its threadgroup memory (64–128 bytes) and
threads per group would allow 96 simdgroups per core, so the register file sets the cap. Halving its accumulators moved the
count only from 109 to 107 and did not lift occupancy: its registers are dominated by weight words, dequantized values,
field temporaries and addresses. Instruments' "Register Pressure Influence" counter read 0 for every kernel recorded,
including these, so that counter does not see this cap.

Registers and residency are measured in \cref{sec:residency-occupancy-measurements}; the declared count is inert for execution and residency
(rule 1 of \cref{sec:residency-occupancy-rules}; MM~\mm{10}, MM~\mm{25.115}).

\textbf{Evidence.} MM~\mm{10}, MM~\mm{11}, MM~\mm{13}, MM~\mm{19}, MM~\mm{25.98.1}, MM~\mm{25.109.1}, MM~\mm{25.111}, MM~\mm{25.115}, MM~\mm{25.144.4}.

\section{Residency and occupancy measurements}\label{sec:residency-occupancy-measurements}

The residency rules (what the hardware guarantees, forward progress) are in \cref{ch:device-execution-model}. Each residency
count here is given with its kernel and conditions, because the counts differ by kernel and must not be merged.

Metal's public API has no occupancy counter; it offers only \code{GPUTimestamp} (measured, MM~\mm{25.48}; \cref{sec:dispatch-encoder-command-buffer}).
\code{maxTotalThreadsPerThreadgroup} reads 1,024 for every pipeline,
from 4 to 512 live scalars per thread and 64 to 32,768 bytes of threadgroup memory, so it does not respond to pressure;
\code{staticThreadgroupMemoryLength} tracks the request (measured, MM~\mm{25.47}, MM~\mm{11}). Residency was therefore measured with probes
built for the purpose, and later read through Instruments counters (\cref{sec:dispatch-encoder-command-buffer}).

An early reading (recon section 167), since corrected, found time flat from 20 to 116
registers per lane and concluded occupancy was "slot-limited, not register-limited". Above the ALU utilisation
ceiling (prior art places it at about 24 simdgroups per core) occupancy can fall a long way with no change in time, so
the flat curve said nothing about occupancy. A direct atomic counter (arrival add, peak max, fixed spin,
exit decrement) then read occupancy falling 44.6 to 32.1 simdgroups per core over 22 to 126 registers. The timings
stand; the occupancy claim was withdrawn (MM~\mm{19}, MM~\mm{20}, MM~\mm{25.48}, MM~\mm{25.49}).

\Cref{tab:residency-counts} gives the counts, each with its kernel.

\begin{xltabular}{\linewidth}{@{}X>{\hsize=0.65\hsize}X>{\hsize=1.77\hsize}X>{\hsize=0.61\hsize}X@{}}
\caption{Resident simdgroups by kernel and probe, with the conditions of each count.}\label{tab:residency-counts}\\
\toprule
Kernel and probe & Conditions & Resident simdgroups & Source \\
\midrule
\endfirsthead
\tablecontinued{4}
\toprule
Kernel and probe & Conditions & Resident simdgroups & Source \\
\midrule
\endhead
Apple-compiled direct atomic counter (recon 170) & 22 to 126 registers & 44.6 falling to 32.1 per core (x1.39); about 40 at 44 registers (44.5 at
30, 39.6 at 46); 32.0 at 98 & measured, MM~\mm{19}, MM~\mm{25.49}, MM~\mm{25.111.2} \\*
Recon 134 kernel, 1,024-thread threadgroups & 40 threadgroups at 126 registers & peak exactly 640 resident, 32 per core; resident falls from about
890–930 at 22 registers to 530–660 at 126 (25–30 percent, no step),
worth 1.00 to 0.95 in throughput capacity & measured, MM~\mm{11} \\*
This compiler's residency counter, straight-line, 12,800
instructions, one simdgroup per threadgroup & warm clock, launched past capacity (S = 4,096), 3 rounds & 95.0 per core at 19 registers used; 71–74 at 33; 67–69 at 49; 66–69 at
65; 65–66 at 81; 64–66 at 97; 68–69 at 113; 64–68 at 125 & measured, MM~\mm{25.121} \\*
Same probe, declared against used & S = 1,024 and 2,048 & declaring 125 while using 19 peaks with using 19 (1,814–1,896 at S =
2,048); warm, using 125 holds 1,024 of 1,024 in 5 of 5 & measured, MM~\mm{25.115} \\*
Threadgroup-memory handle, 128-thread threadgroups, 4,096 threadgroups & 9 to 94 registers & at least 8 simdgroups per core at 94 registers (two threadgroups still
co-resident) & measured, MM~\mm{25.49} \\*
Forward-progress probe, single-simdgroup threadgroups & spinning residents & up to 1,200 threadgroups progress concurrently; first resident wave at
2,000 launched about 1,690 (about 84 per core at that register count);
1,024-thread threadgroups: 200 ran with no consumer starving; at 800,
480 of 799 consumers spun to the poll cap, because the producer,
dispatched last, was beyond residency. That run does not measure how
many were resident at once, since waves can reach the cap one after
another. These are the tested configurations, not a guarantee; the
restriction they motivate (do not depend on preemption; size a waiting
grid to residency) is 
\cref{sec:forward-progress-in-order}'s & measured, MM~\mm{25.141.1} \\*
Instruments, B~8 batched qkv decode kernel & 109 registers & occupancy 25–27 percent of a 96-simdgroup-per-core budget & measured counters, MM~\mm{25.144.4} \\*
\bottomrule
\end{xltabular}

\begin{itemize}
\item The Apple-compiled counter (32.1 at 126) and the recon-134 kernel (32 per core at 126) agree on the floor near 32
  for Apple-compiled kernels at high register counts; MM~\mm{11}'s gradual 25–30 percent fall and the counter's 1.39-fold
  fall describe the same shape on different kernels.
\item This compiler's counter reads about twice as many (64–69 per core from 49 to 125 registers) on a different kernel
  (straight-line code, its own allocator). Why the two kernels differ is open (MM~\mm{25.121}). A scheduler must use the
  floor measured for its own kernel family; \code{agxforge/g17/tensorsched.py} carries both (64 per core from 49 to 125
  registers for this compiler's kernels, 32.1 at 126 for Apple's).
\item Residency follows the registers a program uses, not the count it declares (rule 1 of \cref{sec:residency-occupancy-rules}; MM~\mm{25.115}).
\item \emph{Correction:} MM~\mm{25.115}'s provisional per-core figures are superseded by MM~\mm{25.121} (\cref{app:a}).
\item A fixed-pool model predicts 7.8 simdgroups at 126 registers where 32.1 are measured, so residency is not one pool
  divided by demand; the physical capacity and allocation granularity are open (\cref{sec:residency-occupancy-rules}; MM~\mm{25.49}, MM~\mm{20}).
\end{itemize}

Threadgroup memory also acts as a residency handle (measured, MM~\mm{25.49}). The 57,344-byte per-core budget (rule 4 of
\cref{sec:residency-occupancy-rules}) shows as a step: 28,672 bytes per threadgroup runs 0.00304 s, 28,688 bytes 0.00382 s (+25 percent for 16
bytes). The step persists at 94 registers (ratio 1.047), and its size tracks arithmetic intensity, not registers: a
24-fma control at 10 registers gives 1.041. The 32,768-byte per-threadgroup limit is \cref{ch:memory}'s.

\subsection*{Per-core simdgroup staircase}

\Cref{tab:simdgroup-staircase} gives the median ns per iteration (8 legacy issues) of a retained Apple-compiled legacy
\code{simdgroup\_multiply\_accumulate} loop with 8 independent accumulators and 2,048 trips, at a warm clock (a dispatch of
at least 4.7 ms before every timed one, a fixed 4-simdgroup-per-core reference dispatch after each), grid = 20 ×
simdgroups per core, threadgroup 32 (measured, MM~\mm{25.122}).

\widetable
\begin{longtable}{@{}rr@{}}
\caption{Median time per iteration of the legacy MMA loop against simdgroups per core.}\label{tab:simdgroup-staircase}\\
\toprule
Simdgroups per core & ns per iteration \\
\midrule
\endfirsthead
\tablecontinued{2}
\toprule
Simdgroups per core & ns per iteration \\
\midrule
\endhead
1–8 & 24.5–24.6 \\*
9–12 & 30.7–31.2 \\*
13–17 & 54.6 \\*
18–21 & 60.2–60.9 \\*
22–30 & 83.6–84.3 \\*
31–34 & 85–90.6 \\*
35–42 & 113.4–114.5 \\*
43–45 & 119.6–120.7 \\*
46–48 & 143.0–143.6 \\
\bottomrule
\end{longtable}

\begin{itemize}
\item Rungs at 24.6 + 29.6k ns (k = 0..4: 24.6, 54.0–54.6, 83.7, 113.3, 142.9), each with a +6~ns half-step about 4
  simdgroups before the next full rung (30.6, 60.2, 89.9, 119.5); a full rung spans about 11 simdgroups per core.
\item Models: total work as a line, refuted (steps of 6–24~ns against a 0.5 percent floor); a threadgroup granule, refuted
  (threadgroup sizes 32–128 equal within 1.5 percent at equal work); a residency or wave granule, consistent.
\item The first two boundaries fall exactly per core, at 161 and 241 threadgroups (20 × 8 + 1, 20 × 12 + 1). Later
  boundaries are windows (358–368, 442–460, 613–641, 675–686, 840–855, 891–923 threadgroups) inside which each dispatch
  of the same grid lands on one of two adjacent rungs; the reference reads 24.54~ns either way, so it is not the clock;
  outside the windows one grid repeats within 0.33 percent.
\item The rung counts simdgroups rather than chains in flight. With 2, 4, 8 or 16 accumulators per simdgroup the boundaries
  sit at the same simdgroups per core within one; ns per MMA at 1–8 simdgroups per core is 11.87, 5.99, 3.06 and 1.61,
  a constant about 24~ns per iteration, so up to 16 independent chains hide behind one dependent-chain latency. Time =
  (one simdgroup's iteration time) × f(simdgroups per core), with f about 1 up to 8, 1.25 at 9–12, 2.2 at 13–17, 2.5 at
  18–21.
\item The cause inside the core is open; it is not a residency cap, since the residency counter holds 64–95 per core
  (MM~\mm{25.121}).
\end{itemize}

\emph{Resolution of U3.} A "bimodal legacy rate" that flipped from dispatch to dispatch was first read as a
two-state per-issue rate with a constant 9/8 ratio (recon 158), then as a 79.4~ns quantum (recon 164, in fact the
slope of a line). Both were measurement artefacts: the occupancy axis had been mislabelled by 32 times and
oversubscribed (MM~\mm{25.85}, MM~\mm{25.86}), and runs shorter than about 1.5~ms ran at the idle clock (a v2 sweep with a 0.5~ms
warm-up saw medians of 29.9 and 49.6~ns in the two halves of one run). On the corrected axis at a warm clock the
bimodality is confined to the boundary windows above (MM~\mm{25.122}, MM~\mm{19}).

In a related legacy measurement, one accumulator chained into itself measures latency rather than throughput: 1 chain 129.5
GMAC/s, 2 give 259.1, 4 give 518.1, then 8 and 16 fall to 395.6 and 199.0 as accumulators outgrow the registers
(measured, MM~\mm{25.56}).

\textbf{Evidence.} MM~\mm{11}, MM~\mm{19}, MM~\mm{20}, MM~\mm{25.47}, MM~\mm{25.48}, MM~\mm{25.49}, MM~\mm{25.56}, MM~\mm{25.111.2}, MM~\mm{25.115}, MM~\mm{25.121},
MM~\mm{25.122}, MM~\mm{25.141.1}, MM~\mm{25.144.4}.

\section{Memory bandwidth, latency and instruction fetch}\label{sec:memory-bandwidth-footprint}

Sustained read bandwidth by footprint (measured, MM~\mm{7}, streaming microbenchmark) is 1.3–4.8 TB/s up to 5 MiB
(on-core, reported as a lower bound), 520–528~GB/s at 10–20 MiB (the SLC plateau) and 281–294~GB/s from 40 MiB
(DRAM), with the step between 20 and 40 MiB; independent of threadgroup shape at 1, 8 and 32 simdgroups.

DRAM bandwidth depends on the instrument; \cref{tab:dram-bandwidth} gives the figure the sources report for each.

\begin{xltabular}{\linewidth}{@{}>{\hsize=1.20\hsize}X>{\hsize=1.22\hsize}X>{\hsize=0.58\hsize}X@{}}
\caption{DRAM read bandwidth by measuring instrument.}\label{tab:dram-bandwidth}\\
\toprule
Instrument & Read & Source \\
\midrule
\endfirsthead
\tablecontinued{3}
\toprule
Instrument & Read & Source \\
\midrule
\endhead
streaming microbenchmark by footprint, 40 MiB and above & 281–294~GB/s (row M7: 283–294) & MM~\mm{7}, MM~\mm{20} \\*
1,024-thread threadgroups, 256~MB cold streams & 298.5~GB/s (peak; supersedes the 272~GB/s used before) & MM~\mm{0.15}, MM~\mm{25.141.13} \\*
grid-stride float4, 1,048,576 threads, 256 MiB and 1 GiB, two rounds & 269.6–275.7~GB/s; write 241.0–252.9; copy 243.5–255.6 (read plus write
counted); 64 MiB reads 218–264 & MM~\mm{25.121} \\*
static performance model floor & 280~GB/s (a model constant) & MM~\mm{25.144.12} \\*
\bottomrule
\end{xltabular}

The sources use 298.5~GB/s as the roofline peak, as this chapter does; a simple grid-stride kernel sees about 270. The third instrument agrees with M7 to within about 6 percent (MM
25.121).

\Cref{tab:load-to-use-latency} gives load-to-use latency by footprint (measured, MM~\mm{25.121}; 32 lanes each chasing a
random cycle, warm clock, least-squares slope over five chain lengths).

\begin{xltabular}{\linewidth}{@{}XX@{}}
\caption{Load-to-use latency of a dependent load by footprint.}\label{tab:load-to-use-latency}\\
\toprule
Footprint & ns per dependent load \\
\midrule
\endfirsthead
\tablecontinued{2}
\toprule
Footprint & ns per dependent load \\
\midrule
\endhead
16 KiB (L1) & 30.5 (about 49 cycles) \\*
4 MiB & 195 (192–238) \\*
1 GiB, 16k-64k hops (DRAM) & 405–415 \\*
\bottomrule
\end{xltabular}

The static model uses 30 / 96 / 195 / 313 / 422~ns at 16 KiB / 256 KiB / 4 MiB / 64 MiB / 256 MiB (MM~\mm{25.144.12}); the
256 KiB and 64–256 MiB rows do not fit one line in the probe and are model inputs. Issuing a load costs 9.3~ns in an
unwaited chain (measured, MM~\mm{25.97}). \emph{Correction:} MM~\mm{25.97}'s "about 121~ns" L1 latency was the L1 row at an idle clock
(MM~\mm{25.121}).

Measured cache effects in two real kernels:

\begin{itemize}
\item Decode attention keeps K and V resident in the last-level cache between dispatches up to about 4.2~MB (1,024 keys),
  at 18–20~ns per key; past it about 29~ns per key, about 140~GB/s, under half the DRAM peak (MM~\mm{25.181}).
\item Spare threadgroups warming the cache save up to 10~µs per small-to-big kernel pair (7 percent) when the warm set fits
  the small kernel's duration (4~MB with 256 extra threadgroups: 143.2 to 133.3~µs); decoy loads save nothing; 8~MB
  lengthens the small kernel (MM~\mm{25.141.5}).
\end{itemize}

\subsection*{Instruction fetch in hot loops}

The coarse sweep and the step-location sweep in \cref{tab:instruction-fetch} measure different things.

\begin{xltabular}{\linewidth}{@{}>{\hsize=0.81\hsize}X>{\hsize=1.66\hsize}X>{\hsize=0.53\hsize}X@{}}
\caption{Instruction-fetch cost against loop size, from two sweeps and two follow-up measurements.}\label{tab:instruction-fetch}\\
\toprule
Measurement & Result & Source \\
\midrule
\endfirsthead
\tablecontinued{3}
\toprule
Measurement & Result & Source \\
\midrule
\endhead
Per-instruction time against loop code size, coarse sweep & grows sixfold from 1~KB to 26~KB at one simdgroup per core (0.411 to
2.49~ns), 1.85-fold at 64, then flattens & measured, MM~\mm{7} \\*
Location of the step, seventeen-point sweep at fixed instruction mix & the variable is the loop body: 15,971 bytes fast, 16,595
partial, 17,247 slow, bracketing 16 KiB (1,363 bytes of the
kernel are prologue) & measured, MM~\mm{19} \\*
Cost of crossing it & a loop body over 16 KiB costs 2.81 times on one simdgroup and 1.15 times
at 64 simdgroups per core; cold straight-line code costs 1.91~ns per
byte for one simdgroup & measured, MM~\mm{25.144.12} \\*
At high occupancy & at 26–64 simdgroups of the same code per core fetch is shared and
crossing costs nothing extra: register-O attention bodies of 21.6–22.5
KB folded to 13.1~KB ran at exactly the instruction-count ratio (1,169
against 739~µs = 1.58; 1,829 against 1,141 instructions = 1.60) & measured, MM~\mm{25.144.12} \\*
\bottomrule
\end{xltabular}

The 16 KiB loop-body bracket is where the per-instruction cost steps; the 26~KB figure is where the
coarse sweep stopped rising. The hot-loop guideline "about 26 KiB, roughly 2,800 instructions" (MM~\mm{0.1}, MM~\mm{7}) is
superseded as a design rule: keep a hot loop body under 16 KiB (the fast point is 15,971 bytes). At 26 or more
simdgroups per core the cliff costs nothing measurable; at intermediate occupancy fetch is unmeasured (open). Neither
fetch form improved the static model (\cref{sec:cost-models-their}). An earlier whole-kernel bracket is superseded (\cref{app:a}).

\textbf{Evidence.} MM~\mm{0.15}, MM~\mm{7}, MM~\mm{19}, MM~\mm{20}, MM~\mm{25.97}, MM~\mm{25.121}, MM~\mm{25.141.5}, MM~\mm{25.141.13}, MM~\mm{25.144.12}, MM~\mm{25.181}.

\section{Dispatch costs and counters}\label{sec:dispatch-encoder-command-buffer}

\Cref{tab:dispatch-costs} lists the costs (measured, MM~\mm{0.15}).

\begin{xltabular}{\linewidth}{@{}>{\hsize=0.84\hsize}X>{\hsize=1.33\hsize}X>{\hsize=0.83\hsize}X@{}}
\caption{Measured dispatch, encoder and command-buffer costs.}\label{tab:dispatch-costs}\\
\toprule
Item & Cost & Conditions \\
\midrule
\endfirsthead
\tablecontinued{3}
\toprule
Item & Cost & Conditions \\
\midrule
\endhead
dependent tiny dispatch & 3.7~µs each in a serial encoder; 4.2~µs with a concurrent encoder plus
scoped or resource barriers (MLX's pattern) & chain of one-thread dispatches each reading the previous word (MM
\mm{25.141.10}, \mm{25.141.13}) \\*
extra compute encoder in one command buffer & about 6.5~µs each & 172 encoders per token: 9.275 against 8.163~ms as one encoder (MM
\mm{25.142.1}) \\*
concurrent encoder with scoped buffer barriers & 0.03–0.07~ms per token slower than serial & pipelined q8 token, 196 dispatches (MM~\mm{25.142.1}) \\*
tracked buffers in one serial encoder against untracked buffers ordered
by \code{MTLSharedEvent} & tracked faster by 0.2–0.3~ms per token & ABBA, both precisions (MM~\mm{25.142.1}) \\*
indirect command buffer & direct against indirect 4.82 against 5.83~µs per dispatch (64 × 1
threadgroups), 4.66 against 5.81 (64 × 32); no gain for large dispatches
(143.98 against 145.59~µs; 131.80 against 128.41 at 32~MB) & MM~\mm{25.141.7} \\*
pipeline change & not measurable: one uber pipeline for FFN and w2 behind skip branches
101.9~µs per dispatch shared, 101.8 separate, 99.9 same-pipeline floor & MM~\mm{25.141.2} \\*
grid barrier instead of dispatch boundaries & 133.1 against 120.8~µs per 32~MB phase (fenced, capped barrier among
resident threadgroups; no barrier reached its cap in these runs, 
\cref{sec:forward-progress-in-order}) & MM~\mm{25.141.1} \\*
\bottomrule
\end{xltabular}

\begin{itemize}
\item Pipeline changes should not be counted as a cost (MM~\mm{17} rule 39). A constant-drain fusion model (6.54~µs per pipeline
  change) was refuted: removing a change saved 0.6–2.5~µs, not 6.5; an in-chain difference alone is not a per-switch
  cost, and a fusion saves about the work it removes (measured, MM~\mm{25.142.4}).
\item An in-chain layer A/B predicted the graph at 0.75–1.07 times for changes of 8–22~µs per layer, and small savings
  reached the graph only on the path every token waits for, so changes are judged in the whole graph (MM~\mm{17} rule 41).
\item A prefill GEMM grid with row groups in the high bits of the threadgroup id ranked differently alone and in the graph: 9
  percent faster alone with weights resident and 15 percent slower in the graph, because row groups reading one column
  block were \code{grid\_n} threadgroups apart and each re-streamed it from DRAM (measured, MM~\mm{25.157}).
\end{itemize}

\subsection*{Counters by instrument}

What the sources say about counters differs by instrument, and each statement is correct for the instrument it
describes (\cref{tab:counter-instruments}).

\begin{xltabular}{\linewidth}{@{}>{\hsize=0.81\hsize}X>{\hsize=1.63\hsize}X>{\hsize=0.56\hsize}X@{}}
\caption{Performance counters available through each instrument.}\label{tab:counter-instruments}\\
\toprule
Instrument & What it gives & Source \\
\midrule
\endfirsthead
\tablecontinued{3}
\toprule
Instrument & What it gives & Source \\
\midrule
\endhead
Metal's public counter API & one counter set, \code{timestamp} (\code{GPUTimestamp}), sampled at
stage (encoder) boundaries only: \code{supportsCounterSampling} stage 1,
dispatch 0; sampling at a dispatch boundary asserts & measured, MM~\mm{25.48}, MM~\mm{25.141.6}, MM~\mm{0.15} \\*
Metal System Trace & command buffers and encoders only; the only GPU counter is "RT Unit
Active"; it records the GPU performance state (Minimum, Medium, Maximum)
per interval & measured, MM~\mm{0.15}, MM~\mm{25.142.1} \\*
Instruments \code{xctrace record} with a GUI-saved "Performance
Limiters" template (\code{counterscounterprofile = 13},
\code{g17-limiters.tracetemplate}) & 89 counters, including Kernel Occupancy, Compute SIMD Groups Inflight,
Instruction Throughput Limiter and Utilization, ALU / F32 / F16 /
Integer limiters, L1 Cache Limiter, Buffer L1 Read/Write Bandwidth, GPU
Read/Write Bandwidth, Last Level Cache Bandwidth and Limiter, Memory
Request Stalls Influence, Register Pressure Influence, Occupancy Manager
Target, Neural Accelerator Limiter & measured, MM~\mm{25.144.4} \\*
\bottomrule
\end{xltabular}

The private \code{GPURawCounter.framework} route fails: the sampler-source query (selector 261) returns
kIOReturnNotPrivileged, and root does not help; it needs the entitlement \code{com.apple.private.agx.performance-spi}, and
Apple's \code{MTLReplayer}, which carries it, is killed on launch by AMFI (measured, MM~\mm{25.144.4}). \emph{Correction:} the "M4
status" paragraph inside MM~\mm{25.144.4} predates the template recording later in that section (\cref{app:a}).
"Timestamps only" (MM~\mm{0.15}, MM~\mm{25.48}, MM~\mm{25.141.6}) remains true of Metal's public API.

The Instruments counters are reduced as follows (MM~\mm{25.144.4}). Per-core counters are medians over busy samples (Kernel Occupancy
at least half its own 90th percentile); bandwidths are 90th percentiles from a separate sampler. Absolute "GPU read
bandwidth" is not a DRAM figure when a looped weight stays LLC-resident (MLX's about 500~GB/s exceeds the 298.5~GB/s
DRAM peak). Register Pressure Influence read 0 for every kernel recorded (\cref{sec:registers-spills-occupancy}).

No hardware clock is readable inside a kernel: none of the 67 special-register indices of
\op{14059} is a clock, \code{llvm.readcyclecounter} crashes the compiler service, and a GPU-side
ticker slows five-fold under load. The working clock is the CPU's \code{mach\_absolute\_time} (41.667~ns per tick) written
into 512 lines of an untracked write-combined buffer, read by a one-thread stamp dispatch (about 1~µs each); the
stamped span matches the command buffer to 0.14–0.5 percent (measured, MM~\mm{25.141.13}). One encoder per dispatch is not
a faithful timeline (a layer read 276~µs that way against 330 in one encoder) (measured, MM~\mm{0.15}; MM~\mm{17} rule 42).

Below Metal, the submission interface is \cref{ch:submission-below-metal}'s: a command queue is
created through selectors 7, 16 and 28, and each Submit enters the kernel through \code{IOConnectTrap4} (one record) or
\code{IOConnectCallMethod} selector 0x1d (several records) (\cref{sec:submit-record-call}). The resident runtime's warm Submit-to-callback
split is in \cref{sec:completion-callbacks-their}.

\textbf{Evidence.} MM~\mm{0.15}, MM~\mm{17}, MM~\mm{25.48}, MM~\mm{25.141.1}, MM~\mm{25.141.2}, MM~\mm{25.141.6}, MM~\mm{25.141.7}, MM~\mm{25.141.13},
MM~\mm{25.142.1}, MM~\mm{25.142.4}, MM~\mm{25.144.4}, MM~\mm{25.157}.

\section{GPU clock and power}\label{sec:clocks-power-gpu}

The top clock and the governor are described in \cref{sec:clocks-performance-states}. Clock effects seen in these measurements:

\begin{itemize}
\item Cold runs are slow (the rule, the 1.5~ms idle-clock threshold and the warm-up are \cref{sec:clocks-performance-states}). For example,
  a cold one-threadgroup pointer chase was linear in chain length (206, 405, 802, 1,640~µs at 4 MiB) and
  then fell for longer chains (measured, MM~\mm{25.121}, MM~\mm{25.122}).
\item Short command buffers drop the performance state: mlx-lm's seven per token spent 27–78 percent of decode in
  "Medium" (\cref{sec:clocks-performance-states}; MM~\mm{0.15}).
\item A clock reference does not detect contention for cores. Under another workload (Device Utilization 31 percent) the same FFN
  projection timed 185–194~µs against 145–155 idle while the reference dispatch read 50.3–50.6~µs in both (measured,
  MM~\mm{25.124.6}).
\end{itemize}

\Cref{tab:gpu-rail-power} gives measured power on the GPU rail. Its rows come from two instruments and two engines, so they do not form a like-for-like ratio.

\begin{xltabular}{\linewidth}{@{}>{\hsize=1.48\hsize}X>{\hsize=0.91\hsize}X>{\hsize=0.91\hsize}X>{\hsize=0.70\hsize}Xl@{}}
\caption{GPU-rail power, work rate and efficiency for each arm.}\label{tab:gpu-rail-power}\\
\toprule
Arm & GPU rail & Work & Efficiency & Source \\
\midrule
\endfirsthead
\tablecontinued{5}
\toprule
Arm & GPU rail & Work & Efficiency & Source \\
\midrule
\endhead
idle & 32.3 mW & – & – & MM~\mm{25.56} \\*
fp32 fma & 4,948 mW & 12.4 GMAC/s & 2.5 GMAC/J & MM~\mm{25.56} \\*
legacy engine, \op{838}, four
independent accumulators & 9,973 mW & 184.4 GMAC/s & 18.5 GMAC/J & MM~\mm{25.56} \\*
legacy engine, \op{838}, one dependent
chain & 19,001 mW & 129.6 GMAC/s & 6.8 GMAC/J & MM~\mm{25.56} \\*
idle (second instrument) & 31.6 mW & – & – & MM~\mm{25.60} \\*
fp32 FMA, serially dependent chain & 27,917.5 mW & 0.071 TFLOP/s (latency-bound) & – & MM~\mm{25.60} \\*
TensorOps \code{matmul2d} 32×32×64 half to float & 62,768.4 mW (reproduced at 62,518.0, 0.40 percent) & 4.00 TFLOP/s, 12 percent of 33.5 & 2.06 uJ per 32×32×64 MMA & MM~\mm{25.60} \\*
\bottomrule
\end{xltabular}

\begin{itemize}
\item MM~\mm{25.56}: 5 s windows, three trials, grid $2^{20}$, GPU at 1,620~MHz and 100 percent active residency in every working
  arm. The dependent chain draws 1.9 times the power of the four-chain kernel for less work (trials within 40 mW).
  A kernel stalled on memory behaves differently: the repository README gives 5,607 mW computing against 1,681 mW
  stalled on memory. Its "MMA" rows are the legacy engine, not the tensor unit.
\item MM~\mm{25.60}: filled 6 s windows; liveness ratios (runtime against an eight-fold loop increase) tensor 8.05, ALU 7.82.
  These are operating points, not peaks; "2.25 times" is not the accelerator's power ratio over a saturated ALU.
\item \emph{A dead-code measurement.} The first tensor attempt stored its result under \code{if (lane == 1024)}, the
  compiler deleted the loop, and it read 9.5~W; dispatch time did not move with the loop count (256 and 32,768
  iterations both 0.010~ms), which showed that the loop was gone. Every power figure here carries a liveness ratio for that reason (MM~\mm{25.60}).
\item Killing the host left the GPU at 1,620~MHz, 5,588 mW and 0 percent idle with every process reporting 0 GPU ms/s
  (MM~\mm{17} rule 35).
\end{itemize}

\emph{Correction:} MM~\mm{20} ("Hardware identity") says "Power was not measured"; MM~\mm{25.56} and MM~\mm{25.60} measured it.

\textbf{Evidence.} MM~\mm{0.15}, MM~\mm{8}, MM~\mm{17}, MM~\mm{20}, MM~\mm{25.56}, MM~\mm{25.60}, MM~\mm{25.121}, MM~\mm{25.122}, MM~\mm{25.124.6}.

\section{Cost models}\label{sec:cost-models-their}

The project has two cost models. \code{tensorsched} prices uniform work in measured kernel families and bounds
non-uniform work from below; \code{g17perfmodel} prices any compiled program from its bytes, with a median error of
15 percent, good for attention and poor for GEMM tile shapes and DRAM-bound kernels.

\subsection*{Scheduler model (\code{agxforge/g17/tensorsched.py})}

The model (MM~\mm{16}, MM~\mm{25.124}, and the module docstring) is:

\begin{lstlisting}
T_latency    = max over simdgroups of (steps of that simdgroup) x L_step
T_throughput = (total steps) x C_step / cores
T_memory     = unique bytes / BW(footprint)
time         = max(T_latency, T_throughput, T_memory)
\end{lstlisting}

This is the standard makespan lower bound, tight only when every simdgroup does the same number of steps.

\begin{itemize}
\item \emph{Accuracy on uniform work} (measured, MM~\mm{16}, MM~\mm{25.124}): attention median 2.7 percent (20 of 23 rows within 10
  percent; worst −16.5 percent at 128 × 2,048 keys fp8); weight-stationary stage 1, 5.1 percent stated, 3.3 percent on
  the retained rows (worst −6.0); token-parallel FFN, 3.3 percent stated over 16 rows, of which 7 come from an
  unretained log; the 9 retained rows give 4.7 percent (worst +11.9, bf16 at 10 threadgroups).
\item \emph{Domain limit from a non-uniform workload} (measured, MM~\mm{16}). A causal attention gives query tiles triangular trip
  counts; the model over-predicts the diagonal skip's benefit by 20.9 percent at 256 tasks and 37.1 percent at 1,024,
  and refitting both constants on the same kernels' uniform runs leaves the error. For unequal work its predicted
  speedup is a ceiling and its time a lower bound. The uniform two-regime shape is also too sharp (time rises
  1.46 times for 8 times the tasks, then 3.14 times for the next 4 times; the best fit leaves ±12 percent).
\item \emph{The scheduling term} (model and measured, MM~\mm{25.124.1}). \code{T\_schedule = list\_makespan\_ns}: tasks dealt in threadgroup
  index order round-robin to 20 cores, then one per freed slot; a core with n resident tasks advances each one step per
  \code{max(L\_step, n x C\_step)}. On the causal points (e4m3, head width 128, \code{L\_step} 2,868~ns, \code{C\_step} 327~ns) speedup
  optimism falls from 20.9 to 7.5 percent at 256 tasks and from 37.1 to 26.5 percent at 1,024; it stays a lower bound,
  and \code{speedup} refuses to form a ratio from it.
\item \emph{Constants and domain checks} (model, MM~\mm{25.124}): attention step pairs 5.35~µs / 433~ns (fp8, int8) and 5.0~µs / 381
  ns (bf16); a bandwidth ladder of 3 TB/s to 5 MiB, 0.52 TB/s to 20 MiB and 0.28 TB/s from 40 MiB; the chain fit 62.6 /
  5.8~ns (bf16) and 68.2 / 7.0~ns (fp8, int8) with \code{chain\_ns} refusing any occupancy but 1 and 64 per core; residency
  floors 64 per core from 49 to 125 registers (this compiler's kernels) and 32.1 at 126 (Apple's). Every prediction
  carries checks (kernel family, dtype, head width or layer shape, registers 19–126 with the residency floor holding
  the needed simdgroups, footprint on a measured plateau, token range, uniform work); a failed check yields no number,
  and \code{choose} returns the known-safe schedule with \code{fallback=True}.
\item \emph{Calibration across kernel families} (measured, MM~\mm{25.124.2}, MM~\mm{25.126}): the recon constants chose the wrong FFN
  schedule at fp8 48 and bf16 64 tokens on the retained harness; the \code{layerbench} constants pick the measured-fastest
  at all 11 points (blocked weight-stationary to token-parallel between 32 and 48 e4m3 tokens and between 48 and 64
  bf16; one-pass weight-stationary wins at 48 e4m3, 1.119 against 1.254~ms). Head width scales by the MMA count
  (\code{dh/128} for multiples of 32 from 96 to 256, 12 ratios at −6.4 to +2.3 percent; 144 costs 1.25, measured 1.26–1.32;
  width 64 takes 0.355); the affine per-MMA extrapolation under-predicted width 256 by 16–32 percent and is retired.
\item Simdgroups per threadgroup 1, 2, 4 and 8 predict the same attention time; the model keeps the measured default 8
  (1.455~ms, 4.2 percent behind the best, 2 at 1.396) and excludes 16 (16 threadgroups leave 4 cores idle) (MM
  25.124).
\item It chooses schedules but does not lower them, and the production runtime cannot express most schedules it chooses between
  (module docstring).
\end{itemize}

\subsection*{Static model from compiled bytes (\code{tools/g17perfmodel.py})}

The figures in this subsection are model and measured values from MM~\mm{25.144.12}.

\begin{itemize}
\item \emph{Mechanics.} One simdgroup is simulated in order over decoded bytes: each instruction issues one 3.105~ns interval
  after the last and not before its operands; loads publish scoreboard slots and consumers stall until they land
  (tensor fragment loads, \op{12674} and \op{12709}, take the working-set
  latency; \op{12682} the 16 KiB latency); a 16-bit MMA, \op{5106}, result is ready after a dependent-MMA latency;
  loops are walked 4 trips and extrapolated (latches \code{cnt < imm} and \code{p > 0} only). Threadgroups are dealt to 20 cores
  in index order; residency comes from the program's highest register (the MM~\mm{25.121} table) and 3,072 threads per
  core; per-core rate is an exact mean-value analysis of three stations (issue slots, load path, tensor unit) fed into
  the MM~\mm{25.124.1} list schedule; a 280~GB/s DRAM floor; a fixed per-dispatch cost per timing protocol.
\item \emph{Nine fitted constants} (bounded least squares on log error; the tensor unit may not beat 4.95~ns per MMA per core):
  12.5 issue slots per interval per core; 3.43~ns per fragment load; a scalar load at 0.43 fragment loads; 6.83~ns per
  MMA (72 percent of peak); load latency x0.50 of the pointer chase; dependent MMA latency 31.6~ns; MMA issue 1.0
  interval; per-dispatch cost 1.4 / 1.0 / 1.0~µs. The latency and issue constants sit at their lower bounds: they are
  not identifiable from these timings.
\item \emph{Accuracy} over 108 recorded timings (45 GEMM, 54 attention, 9 graph; 16 set aside as noisy): leave-one-table-out
  median 14.8 percent (47 of 92 within ±15 percent); fixed split 15.5 percent; in-sample 9.6 percent; graph
  attention 8.3 percent; hot attention at M ≥ 512 9.9 percent; warm GEMM 19.6 percent; noisy rows 35 percent. The
  ±15 percent target is met in the median and for attention, not for GEMMs.
\item \emph{Where it fails.} GEMM tile shapes do not extrapolate (a 2×2 tile up to 40 percent slow, tall tiles 16–37 percent
  fast; operand sharing among a threadgroup's simdgroups is invisible to it); 1,024-token prefill projections predicted
  144~ms against 113 measured (+27 percent); SwiGLU 9.0 against 7.2~ms (DRAM-bound, and the model has one 280~GB/s
  plateau); per-simdgroup latency is not identifiable. It refuses to run without the opcode class table, because an
  UNKNOWN fallback moved 81 of 108 predictions.
\item \emph{Fetch.} Two fetch terms built from the recorded cliff (1.91~ns per byte cold, 2.81 times past 16 KiB on one
  simdgroup, 1.15 times at 64) both worsened held-out error (23.2 and 17.5 percent against 14.8), so fetch is not in
  the model of record.
\item \emph{Predictions it made} (prediction, then built): SwiGLU fused into w3's GEMM 7.6~ms (built); hardware exp2 at 512-row
  chunks 1.5~ms (built; 679 to 490 instructions a trip); K and V shared across 4 simdgroups through threadgroup memory
  −2.8~ms (slower).
\end{itemize}

\textbf{Evidence.} MM~\mm{16}, MM~\mm{25.124}, MM~\mm{25.124.1}, MM~\mm{25.124.2}, MM~\mm{25.126}, MM~\mm{25.144.12}; \code{agxforge/g17/tensorsched.py}.

\section{Limiting resources by workload}\label{sec:bounds-real-workloads}

The workload is InternLM2.5-1.8B at 4 bits unless stated. Throughput against mlx-lm for each workload is in the
machine model (MM~\mm{25.190}; the headline comparison is in \href{https://github.com/sbryngelson/AGXForge/blob/main/docs/showcase.md}{the showcase}, section 2); the evidence for which resource bounds
each workload is kept here.

\subsection*{Single-stream decode}

Wide projections are memory-bound and narrow ones issue-bound.

\begin{itemize}
\item In the plan (model and measured, MM~\mm{25.124.5}), one bf16 decode step (d 2048, 16 heads x 128, KV 256, one query row)
  has a memory floor of 0.487~ms (136.3~MB), 0.247~ms with fp8 weights.
\item A one-token projection with too few threadgroups is bound by its K chain (measured, MM~\mm{25.124.6}): M 16, N 128, K
  2048 times the same with cold and cache-resident weights (within 1 percent), at 6.5–9.9~GB/s ours and 7.5–16.7~GB/s
  Apple's, against 270. With enough threadgroups the column grid reaches the floor: N 8192 × K 2048 (33.5~MB) on 64
  threadgroups of 128 columns takes 143–145~µs = 231–234~GB/s, 1.2 times its 0.120~ms floor, equal to Apple's
  \code{matmul2d} (142.8–145.2~µs). Split-K on N 2048, K 2048: G 1 84.2, 2 52.1, 4 41.1, 8 42.0, 16 43.3~µs; the plateau
  (about 200~GB/s, cached equal to cold) is not DRAM.
\item \Cref{tab:decode-time-breakdown} shows where 4.73~ms per token goes at a 1K context (measured, MM~\mm{25.181}).

  \begin{xltabular}{\linewidth}{@{}>{\hsize=1.20\hsize}Xl>{\hsize=0.80\hsize}Xl@{}}
\caption{Single-stream decode time per token at a 1K context by part, with bytes streamed and effective bandwidth.}\label{tab:decode-time-breakdown}\\
\toprule
Part & ms & Bytes per token & Effective GB/s \\
\midrule
\endfirsthead
\tablecontinued{4}
\toprule
Part & ms & Bytes per token & Effective GB/s \\
\midrule
\endhead
  FFN (w1, w3, SwiGLU) & 1.76 & 453~MB & 257 \\*
  w2 + residual & 0.90 & 227~MB & 252 \\*
  qkv & 0.48 & 113~MB & 236 \\*
  wo + residual & 0.30 & 57~MB & 188 \\*
  head & 0.36 & 107~MB & 296 \\*
  attention & 0.74 & 107~MB of KV & 144 \\*
  norms and the rest & 0.18 &  &  \\*
  \bottomrule
\end{xltabular}

  Only the head streams at the DRAM peak.
\item Counters on the decode qmv kernels (measured, MM~\mm{25.144.4}): the q4 w2 kernel is issue-bound (instruction
  throughput limiter 77.2 percent, zero memory-stall influence; pinned-weight floor 32.9~µs above the 29.8~µs DRAM
  floor); q8 w2 is memory-bound (instruction limiter 46.7 percent, memory-stall influence 26.6 percent, LLC limiter
  100; 89 percent of the 298.5~GB/s peak); MLX's qmv is latency- and memory-bound at 16.7 percent occupancy (64-thread
  threadgroups), 37–59 percent memory-stall influence.
\item In q4 w2, compute and stream barely overlap (measured, MM~\mm{25.141.15}, at the 4-rows-per-simdgroup shape): real 46.2~µs; x index cache-resident 40.9;
  weight index resident 29.5; both 28.5~µs (the instruction floor); streaming 8.9~MB at 298.5~GB/s needs about 30~µs.
  The real time is close to the sum.
\item The single-stream q4 MLP is load-bound (measured, bounded negative, MM~\mm{25.200}): removing 26 percent of its
  instructions (FFN 393 to 291 instructions, 10 loads either way) moves nothing (186.4–192.1 against 177.9–193.4 tok/s
  at 1,792 tokens).
\item Context costs about 20~ns per key-layer (measured, MM~\mm{25.199}); long-context attention variants (\code{keyblock},
  \code{tgsplit}) do not help in the graph at 1,792 keys.
\item Apple-compiled twins of this project's decode kernels, bit-identical in output, give 10–15 percent higher median
  throughput, so this project's instruction-level code generation costs that much (measured, MM~\mm{25.211}; showcase
  section 2).
\end{itemize}

\subsection*{Batched decode}

The batched projections are occupancy-bound, then fp32-pipe-bound after split-K (all figures measured).

\begin{itemize}
\item The tensor-unit projection (\code{g17qsm}: $Y^{T} = W^{T} X^{T}$, W dequantized once into registers as A, 16 sequences as B)
  costs 446~µs per layer for 16 tokens (27.9~µs a token) against batched qmv's 126 at B~8: 4.5 times per token; about
  76~GB/s, latency-bound per simdgroup (about 700 dependent scalar instructions per trip before 8 MMAs) (MM~\mm{25.166}).
\item Split-K to about 512 threadgroups is the knee at every role (qkv split 4: 39.8~µs at 512 threadgroups; wo split 8:
  23.2; w1/w3 split 2: 76.4; w2 split 8: 75.8) (MM~\mm{25.185}).
\item At B~16 the projections are fp32-pipe-bound, not memory-bound: F32 Limiter 100, F32 Utilization 99.9, instruction
  throughput limiter 78 (p90 98), kernel occupancy 38 (p90 50), memory request stalls 0, GPU read bandwidth p90 122
  GB/s, Neural Accelerator utilization p90 25, F16 utilization 0. Projections are 7.7~ms of a 10.98~ms step (70
  percent) for 0.95~GB of q4, about 2.6 times the DRAM floor. Ablation on w1: base 76.3, no dequant 59.6, no weight
  load 60.9, neither 35.1~µs (MM~\mm{25.194}).
\item Moving the dequant to the fp16 pipe with \op{798}, two per weight, relieves the fp32 pipe:
  w1/w3 74.6 to 65.7~µs (−12 percent), w2 73.9 to 66.6, qkv 39.6 to 35.4, wo 22.1 to 20.0 (MM~\mm{25.196}).
\item At B~8 the w2 projection hits the LLC limiter (100 percent) for this compiler and for MLX's \code{affine\_qmv\_wide}
  alike, so both are held by the same memory limit. The B~8 qkv is register-capped at 25–27 percent occupancy (\cref{sec:registers-spills-occupancy}) (MM~\mm{25.144.4}).
\item MLX's own batched decode is compute-bound past B~2 (quantized) and B~8 (bf16); every precision converges on 19–21~ms
  per step at B~16 (q4 841 tok/s, bf16 750): dequantization costs more than the bytes saved (MM~\mm{25.142.5}).
\item Which mlx-lm batched path is the baseline depends on the model (MM~\mm{25.193}).
\end{itemize}

\subsection*{Prefill}

Every figure in this subsection is measured.

\begin{itemize}
\item \emph{A corrected GEMM measurement.} An earlier comparison put this compiler's prefill GEMMs at 60–69 percent
  of peak against MLX's 82–88. The per-dispatch profile used there did not close (181~ms of GPU time in a 150~ms
  prefill), and the isolated GEMM bench runs about four times the in-graph per-call time. With a closure check
  (dispatch-alone sum 135.8~ms against chained step spans 137.1 and in-process medians 139.6, within 1 percent) the
  in-graph GEMMs are at 80–89 percent of 33.1 TFLOP/s: qkv 256 × 4096 × 2048 148~µs (88 percent), wo 81~µs (80), w1
  290~µs (89), w2 292~µs (89). Pinning the weights cache-resident changed per-role time by 0 to ±5 percent, so weight
  traffic is not the GEMM cost (MM~\mm{25.160}).
\item Prefill attention is latency-bound on the per-block scratch round trip: 32-key blocks cut instructions per 16 keys
  from 490 to about 390 (−20 percent) for −3 percent time (MM~\mm{25.162}). Keeping S and P in registers made attention 14.7
  to 12.1~ms (exact) or 9.6~ms (hardware exp2) (MM~\mm{25.163}).
\item The SwiGLU kernel was bandwidth-bound: 10 bytes an element over 201 M elements is 6.7~ms at 298.5~GB/s, exactly its
  6.74~ms; an fp16 intermediate (6 bytes an element) made it 5.54~ms and arithmetic-bound (MM~\mm{25.183}).
\end{itemize}

\Cref{tab:binding-resources} summarizes what binds each workload.

\begin{xltabular}{\linewidth}{@{}X>{\hsize=1.34\hsize}X>{\hsize=0.69\hsize}X@{}}
\caption{Binding resource for each studied workload.}\label{tab:binding-resources}\\
\toprule
Workload & Binding resource in the studied kernels & Evidence \\
\midrule
\endfirsthead
\tablecontinued{3}
\toprule
Workload & Binding resource in the studied kernels & Evidence \\
\midrule
\endhead
fused attention, FFN layer kernels & instruction issue and ALU (about 46 instructions per MMA) & MM~\mm{8}, MM~\mm{16} \\*
single-stream decode, q8 projections and the head & DRAM bandwidth & MM~\mm{25.144.4}, MM~\mm{25.181} \\*
single-stream decode, q4 w2 & instruction issue (counter) yet insensitive to removing dequant work
(load-bound) & MM~\mm{25.144.4}, MM~\mm{25.200} \\*
one-token projection with few threadgroups & one K chain's latency & MM~\mm{25.124.6} \\*
decode attention past 4.2~MB of KV & LLC capacity, then about 140~GB/s & MM~\mm{25.181} \\*
batched decode B~16 projections & occupancy first (about 256 one-simdgroup threadgroups a projection);
after split-K to about 512, the fp32 pipe, relieved by moving the
dequant to the fp16 pipe & MM~\mm{25.185}, MM~\mm{25.194}, MM~\mm{25.196} \\*
batched decode B~8 w2 & LLC bandwidth (both stacks) & MM~\mm{25.144.4} \\*
prefill GEMMs & tensor and issue, at 80–89 percent of peak & MM~\mm{25.160} \\*
prefill attention & latency of a scratch round trip, until S and P stay in registers & MM~\mm{25.162}, MM~\mm{25.163} \\*
\bottomrule
\end{xltabular}

\textbf{Not known.} The cause of the per-core staircase (MM~\mm{25.122}); why this compiler's residency counter reads
about twice Apple's kernel's (MM~\mm{25.121}); fetch cost at intermediate occupancy (MM~\mm{25.144.12}); physical register
capacity per core (MM~\mm{20}); why Apple-compiled twins are 10–15 percent faster (MM~\mm{25.211}). All figures are from one machine, and the
Metal-path evidence files behind them do not record the OS build (\cref{sec:measured-part-software}).

\textbf{Evidence.} MM~\mm{8}, MM~\mm{16}, MM~\mm{25.124.5}, MM~\mm{25.124.6}, MM~\mm{25.141.15}, MM~\mm{25.142.5}, MM~\mm{25.144.4}, MM~\mm{25.160}, MM~\mm{25.162},
MM~\mm{25.163}, MM~\mm{25.166}, MM~\mm{25.181}, MM~\mm{25.183}, MM~\mm{25.185}, MM~\mm{25.190}, MM~\mm{25.193}, MM~\mm{25.194}, MM~\mm{25.196}, MM~\mm{25.199},
MM~\mm{25.200}, MM~\mm{25.211}.
